% !TEX root = ../Abschlussarbeit_Jan_Höndl.tex
\usepackage{amsmath}
\usepackage{amssymb}
\usepackage{graphicx}
\usepackage{hyperref} 
\hypersetup{
    colorlinks=true,
    linkcolor=black,    % Interne Links (Inhaltsverzeichnis etc.)
    citecolor=black,    % Zitate
    filecolor=black,
    urlcolor=black      % Web-Links
}
\usepackage{subcaption}
\captionsetup{
    font=small,              % Etwas kleiner als der Haupttext
    labelfont=bf,            % "Abbildung 1:" in fett
    labelsep=colon,          % Doppelpunkt nach der Nummer
    format=plain,            % Text bündig unter dem Label
    margin=10pt              % Einrückung links/rechts zur besseren Trennung vom Text
}

% Spezielle Konfiguration für die kleinen Unterabbildungen (falls nötig)
\captionsetup[sub]{
    font=footnotesize,       % Noch ein Stück kleiner für a), b), c)
    labelfont=md             % Label hier nicht fett, sondern normal
}
\usepackage{here}
\usepackage{color}
\usepackage{textcomp}
\usepackage{pdfpages}
\usepackage{multirow}
\usepackage{microtype}


%Zeilenabstände und Seitenränder
\usepackage{setspace}
\singlespacing
\usepackage{geometry}
\geometry{
  a4paper,
  left=40mm,           % Linker Rand (Bindung)
  right=20mm,          % Rechter Rand
  top=30mm,            % Oberer Rand
  bottom=30mm,         % Unterer Rand
  headsep=12.5mm,      % Abstand zu Kopfzeile (1,25 cm)
  footskip=12.5mm,     % Abstand zu Fußzeile (1,25 cm)
  includeheadfoot=false % Ränder gelten ab Blattkante
}

\graphicspath{{Bilder/}}
% Ebene 1: 16 Pkt, Fett, Linksbündig
\setkomafont{section}{\fontsize{16pt}{20pt}\selectfont\bfseries\rmfamily}

% Ebene 2: 14 Pkt, Fett, Linksbündig
\setkomafont{subsection}{\fontsize{14pt}{18pt}\selectfont\bfseries\rmfamily}

% Ebene 3: 12 Pkt, Fett, Linksbündig
\setkomafont{subsubsection}{\fontsize{12pt}{14pt}\selectfont\bfseries\rmfamily}

% Ebene 4: 12 Pkt, Kursiv, Linksbündig
% (In scrartcl ist die 4. Ebene meist \paragraph)
\setkomafont{paragraph}{\fontsize{12pt}{14pt}\selectfont\itshape}

\usepackage{scrlayer-scrpage}
\usepackage{enumitem}
% Optional: Global alle Listen kompakter machen
\setlist[itemize]{noitemsep, topsep=0pt}

\usepackage[
  backend=biber,       % Nutze Biber als Sortier-Engine (wichtig!)
  style=authoryear,    % Zitierstil (z. B. Autor-Jahr-Stil: (Feth 2023))
  % style=numeric,     % Alternativ: Numerischer Stil: [1]
]{biblatex}
\addbibresource{Bibliographie.bib}
\let\cite\parencite    % Ermöglicht geklammerte Zitate mit \cite wie bei natbib
\usepackage[ngerman]{todonotes}

%Abkürzungsbibliothek
\usepackage[printonlyused]{acronym}
%Glossar
\usepackage[automake]{glossaries}
\makeglossaries

\usepackage{fontspec}
\setmainfont{Times New Roman} % Legt Times New Roman als Hauptschriftart fest
%Code
%Siehe auch: http://ctan.ebinger.cc/tex-archive/macros/latex/contrib/listings/lstdrvrs.pdf
\usepackage{listings}
\usepackage{xcolor}
\definecolor{neo4jSpace}{HTML}{FF9F9F}       % Rosa für IfcSpace
\definecolor{neo4jStorey}{HTML}{B3A2C7}      % Lila für IfcBuildingStorey
\definecolor{neo4jBg}{HTML}{12161A}          % Dunkler Graph-Hintergrund
\definecolor{neo4jSelection}{HTML}{A5F3FC}
% Definition der Inline-Boxen für die Labels (ahmt die Neo4j-GUI nach)
%\usepackage{uarial}
\definecolor{json_key}{rgb}{0,0,0.6}
\definecolor{json_string}{rgb}{0,0.5,0}
\lstdefinelanguage{cypher}{
    morekeywords={MATCH, WHERE, RETURN, AND, NOT, OR, WITH, UNWIND, CREATE, MERGE, DELETE, SET, REMOVE},
    sensitive=true,
    morecomment=[l]{//},
    morestring=[b]",
    morestring=[b]',
}

\lstdefinelanguage{IFC}{
    morekeywords={ISO, HEADER, FILE_DESCRIPTION, FILE_NAME, FILE_SCHEMA, ENDSEC, DATA, END, IFCPROJECT, IFCOWNERHISTORY, IFCPERSONANDORGANIZATION, IFCPERSON, IFCORGANIZATION, IFCAPPLICATION, IFCUNITASSIGNMENT, IFCSIUNIT, IFCGEOMETRICREPRESENTATIONCONTEXT, IFCAXIS2PLACEMENT3D, IFCCARTESIANPOINT, IFCSITE, IFCBUILDING, IFCPOSTALADDRESS, IFCBUILDINGSTOREY, IFCWALLSTANDARDCASE, IFCLOCALPLACEMENT, IFCPRODUCTDEFINITIONSHAPE, SHAPEREPRESENTATION, EXTRUDEDAREASOLID, IFCRECTANGLEPROFILEDEF, IFCDIRECTION, IFCRELCONTAINEDINSPATIALSTRUCTURE},
    sensitive=true,
    morecomment=[s]{/*}{*/},
    morestring=[b]'
}

\lstset{
    basicstyle=\ttfamily\small,
    breaklines=true,
    numbers=left,
    numberstyle=\tiny\color{gray},
    keywordstyle=\color{blue}\bfseries,
    stringstyle=\color{orange},
    commentstyle=\color{green!60!black},
    frame=lines,
    captionpos=b
}
\lstdefinelanguage{GraphQL}{
  keywords={interface, implements, type, schema, scalar, input},
  keywordstyle=\color{blue}\bfseries,
  ndkeywords={ID, String, Int, Float, Boolean},
  ndkeywordstyle=\color{purple}\bfseries,
  identifierstyle=\color{black},
  sensitive=true,
  comment=[l]{\#},
  commentstyle=\color{gray}\ttfamily,
  stringstyle=\color{red}\ttfamily,
  morestring=[b]",
}
\lstdefinelanguage{json}{
    basicstyle=\normalfont\ttfamily\small,
    numbers=left,
    numberstyle=\scriptsize,
    stepnumber=1,
    numbersep=8pt,
    showstringspaces=false,
    breaklines=true,
    frame=lines,
    string=[s]{"}{"},
    comment=[l]{:\ },
    morecomment=[s]{"{}}{"},
}

% Globale Einstellungen für den Code-Block
\lstset{
  language=GraphQL,
  backgroundcolor=\color{gray!5},
  basicstyle=\small\ttfamily,
  breaklines=true,
  captionpos=b,
  frame=single,
  numbers=left,
  numberstyle=\tiny\color{gray},
  tabsize=2,
  showstringspaces=false
}


\usepackage{chngcntr}
\counterwithin{figure}{section}
\usepackage[perpage]{footmisc}
\usepackage[ngerman]{babel}

